\section{Introduction}
In \textit{Informational Energetics: A Universal Architecture of Persistence} \cite{meyer_informational_2026}, we introduced Informational Energetics (IE) and established that any persistent system must resist entropic decay by satisfying six architectural pillars derived from control theory, information theory, and thermodynamics: Capacity (the vessel), Map (the model), Protocol (the interface), Governor (the constraint), Toll (the temporal cost), and Margin (the buffer). In \textit{Informational Energetics: Entropic Action} \cite{meyer_ieaction_2026}, we formalized the Entropic Action ($S_\Phi$) which defines the total thermodynamic cost for a persistent system to maintain structural coherence.

While previous work established what a persistent system must \textit{be} and how it must \textit{behave}, neither addressed how it can continue to persist when its regulatory complexity exceeds the capacity of a single control loop. Sub-partitioning can be found in every domain: a biological cell that must simultaneously maintain membrane integrity, metabolize, and replicate partitions into organelles; a corporation partitions into departments. In \cref{sec:nested_persistence}, we define three fundamental rules of \textbf{Nested Persistence} governing this partitioning: \textit{Capacity Partitioning}, \textit{Reversible Encoding}, and \textit{Impedance Matching}. Together, they define how a system partitions into specialized subsystems to ensure continued persistence.

\subsection{The Test: The Vacuum}

The vacuum is the ideal test case. Macroscopic systems frequently utilize external energy for partitioning, but a closed system with no external environment to offset its dissipation is forced to solve the partitioning problem with zero error. It therefore provides the most stringent possible test of whether these rules are universal. To test these rules, we apply them to the $E_8$ substrate derived in \cite{meyer_informational_2026}. This substrate contains four key invariants: the temporal resonance $\Delta=43$, the chiral capacity $\nu=16$, the interaction rank $\sigma=5$, and the topological boundary $\chi=2$, all derived from the $E_8 \rightarrow D_4 \oplus D_4$ projection.

While we derived the functional operator forms in ``Informational Energetics: Entropic Action''\cite{meyer_ieaction_2026}, the numerical coefficients remained undetermined. Standard QFT treats these constants as independent empirical inputs, each determined by measurement. In the IE framework, they emerge as coupled partition coefficients of a single finite capacity budget. Unlike $\alpha^{-1}$, which represents the unpartitioned baseline impedance derived directly from the substrate, the remaining parameters are not fundamental. They emerge strictly as the partition coefficients of the vacuum's finite channel capacity ($\nu = 16$): the unique allocation required for nested subsystems to satisfy global solvency. This paper derives these coefficients. Every value is a rigid, parameter-free prediction: a single disagreement with measurement beyond experimental uncertainty falsifies the $E_8$-Persistence theory.

\subsection{The Geometric Cascade}

The $E_8$ substrate possesses a strictly finite channel capacity ($\nu=16$) and a discrete temporal resonance 
($\Delta=43$). To simultaneously govern gauge interactions, electroweak symmetry breaking, and gravitational geometry, the vacuum must partition this finite bandwidth into nested subsystems. The fundamental constants emerge as the exact partition coefficients of this allocation.

We trace this \textbf{Geometric Cascade} across the three sequential systems of the framework, continuing the numbering established in the companion papers (System 0: architectural requirements, System 1: the substrate, System 2: geometric impedance, and System 3: Entropic Lagrangian dynamics):


\begin{itemize}
    \item \textbf{The Capacity Partition (System 4)}: The $\nu = 16$ chiral capacity is subdivided into non-overlapping geometric allocations. We derive the gauge and flavor couplings: Strong Coupling ($\alpha_s$), Weak Mixing Angle ($\sin^2\theta_W$), Cabibbo Angle ($\theta_C$), and Jarlskog Invariant as exact partition coefficients. Because the temporal and color sectors are geometrically isolated by this partition, we show that the Strong CP problem is structurally prohibited without requiring an axion.

    \item \textbf{The Resonant Interface (System 5)}: A nested subsystem stabilizing local topological boundaries against vacuum resonance is identified as the Higgs mechanism. Global impedance matching dynamically generates the electroweak mass scale: 
    the Higgs VEV ($v$), Fermi constant, $W$-boson mass, Higgs mass, and Electron Yukawa ($y_e$). We derive a value for $M_W$ that lies between the Standard Model global fit and the CDF II measurement.

    \item \textbf{The Geometric Depth (System 6)}: A nested subsystem stabilizing the macroscopic lattice volume is identified as Gravity. Signal attenuation across the geometric bulk derives the Planck scale ($M_P$) and Cosmological Constant ($\Lambda$), which we show accounts for the observed Gravitational Hierarchy and Vacuum Catastrophe.
\end{itemize}

\subsection{Note on Absolute Scale}
A discrete lattice inherently produces pure dimensionless ratios (e.g., $v/m_e \approx 4.8 \times 10^5$). To map these geometric ratios to physical units, \textit{exactly one dimensionful empirical input} is required to define the base unit. Here, we utilize the electron mass ($m_e \approx 0.511$ MeV) as the empirical anchor for the Information-Theoretic Noise Floor; all subsequent mass scales ($v, M_W, m_H, M_P$) follow from this single anchor with zero additional tunable parameters.

This derivation cascade spans 64 orders of magnitude, from the gravitational coupling ($10^{-45}$) to the Planck mass ($10^{19}$ GeV) with \textbf{zero tunable parameters}. Every mass, angle, and coupling flows strictly from the $E_8$-Persistence theory substrate.